% Options for packages loaded elsewhere
\PassOptionsToPackage{unicode}{hyperref}
\PassOptionsToPackage{hyphens}{url}
\PassOptionsToPackage{space}{xeCJK}
%
\documentclass[
  12pt,
  UTF8]{ctexart}
\usepackage{amsmath,amssymb,mathtools,setspace}
\setstretch{1.35}
\usepackage{iftex}
\ifPDFTeX
  \usepackage[T1]{fontenc}
  \usepackage[utf8]{inputenc}
  \usepackage{textcomp} % provide euro and other symbols
\else % if luatex or xetex
  \usepackage{unicode-math} % this also loads fontspec
  \defaultfontfeatures{Scale=MatchLowercase}
  \defaultfontfeatures[\rmfamily]{Ligatures=TeX,Scale=1}
\fi
\usepackage{lmodern}
\ifPDFTeX\else
  % xetex/luatex font selection
  \setmainfont[]{Tinos}
  \ifXeTeX
    \usepackage{xeCJK}
    \setCJKmainfont[]{Noto Serif CJK SC}
          \fi
  \ifLuaTeX
    \usepackage[]{luatexja-fontspec}
    \setmainjfont[]{Noto Serif CJK SC}
  \fi
\fi
% Use upquote if available, for straight quotes in verbatim environments
\IfFileExists{upquote.sty}{\usepackage{upquote}}{}
\IfFileExists{microtype.sty}{% use microtype if available
  \usepackage[]{microtype}
  \UseMicrotypeSet[protrusion]{basicmath} % disable protrusion for tt fonts
}{}
\makeatletter
\@ifundefined{KOMAClassName}{% if non-KOMA class
  \IfFileExists{parskip.sty}{%
    \usepackage{parskip}
  }{% else
    \setlength{\parindent}{0pt}
    \setlength{\parskip}{6pt plus 2pt minus 1pt}}
}{% if KOMA class
  \KOMAoptions{parskip=half}}
\makeatother
\usepackage{xcolor}
\usepackage[a4paper,margin=30mm]{geometry}
\usepackage{color}
\usepackage{fancyvrb}
\newcommand{\VerbBar}{|}
\newcommand{\VERB}{\Verb[commandchars=\\\{\}]}
\DefineVerbatimEnvironment{Highlighting}{Verbatim}{commandchars=\\\{\}}
% Add ',fontsize=\small' for more characters per line
\newenvironment{Shaded}{}{}
\newcommand{\AlertTok}[1]{\textcolor[rgb]{1.00,0.00,0.00}{\textbf{#1}}}
\newcommand{\AnnotationTok}[1]{\textcolor[rgb]{0.38,0.63,0.69}{\textbf{\textit{#1}}}}
\newcommand{\AttributeTok}[1]{\textcolor[rgb]{0.49,0.56,0.16}{#1}}
\newcommand{\BaseNTok}[1]{\textcolor[rgb]{0.25,0.63,0.44}{#1}}
\newcommand{\BuiltInTok}[1]{\textcolor[rgb]{0.00,0.50,0.00}{#1}}
\newcommand{\CharTok}[1]{\textcolor[rgb]{0.25,0.44,0.63}{#1}}
\newcommand{\CommentTok}[1]{\textcolor[rgb]{0.38,0.63,0.69}{\textit{#1}}}
\newcommand{\CommentVarTok}[1]{\textcolor[rgb]{0.38,0.63,0.69}{\textbf{\textit{#1}}}}
\newcommand{\ConstantTok}[1]{\textcolor[rgb]{0.53,0.00,0.00}{#1}}
\newcommand{\ControlFlowTok}[1]{\textcolor[rgb]{0.00,0.44,0.13}{\textbf{#1}}}
\newcommand{\DataTypeTok}[1]{\textcolor[rgb]{0.56,0.13,0.00}{#1}}
\newcommand{\DecValTok}[1]{\textcolor[rgb]{0.25,0.63,0.44}{#1}}
\newcommand{\DocumentationTok}[1]{\textcolor[rgb]{0.73,0.13,0.13}{\textit{#1}}}
\newcommand{\ErrorTok}[1]{\textcolor[rgb]{1.00,0.00,0.00}{\textbf{#1}}}
\newcommand{\ExtensionTok}[1]{#1}
\newcommand{\FloatTok}[1]{\textcolor[rgb]{0.25,0.63,0.44}{#1}}
\newcommand{\FunctionTok}[1]{\textcolor[rgb]{0.02,0.16,0.49}{#1}}
\newcommand{\ImportTok}[1]{\textcolor[rgb]{0.00,0.50,0.00}{\textbf{#1}}}
\newcommand{\InformationTok}[1]{\textcolor[rgb]{0.38,0.63,0.69}{\textbf{\textit{#1}}}}
\newcommand{\KeywordTok}[1]{\textcolor[rgb]{0.00,0.44,0.13}{\textbf{#1}}}
\newcommand{\NormalTok}[1]{#1}
\newcommand{\OperatorTok}[1]{\textcolor[rgb]{0.40,0.40,0.40}{#1}}
\newcommand{\OtherTok}[1]{\textcolor[rgb]{0.00,0.44,0.13}{#1}}
\newcommand{\PreprocessorTok}[1]{\textcolor[rgb]{0.74,0.48,0.00}{#1}}
\newcommand{\RegionMarkerTok}[1]{#1}
\newcommand{\SpecialCharTok}[1]{\textcolor[rgb]{0.25,0.44,0.63}{#1}}
\newcommand{\SpecialStringTok}[1]{\textcolor[rgb]{0.73,0.40,0.53}{#1}}
\newcommand{\StringTok}[1]{\textcolor[rgb]{0.25,0.44,0.63}{#1}}
\newcommand{\VariableTok}[1]{\textcolor[rgb]{0.10,0.09,0.49}{#1}}
\newcommand{\VerbatimStringTok}[1]{\textcolor[rgb]{0.25,0.44,0.63}{#1}}
\newcommand{\WarningTok}[1]{\textcolor[rgb]{0.38,0.63,0.69}{\textbf{\textit{#1}}}}
\setlength{\emergencystretch}{3em} % prevent overfull lines
\providecommand{\tightlist}{%
  \setlength{\itemsep}{0pt}\setlength{\parskip}{0pt}}
\setcounter{secnumdepth}{-\maxdimen} % remove section numbering
\ifLuaTeX
\usepackage[bidi=basic]{babel}
\else
\usepackage[bidi=default]{babel}
\fi
\babelprovide[main,import]{chinese}
\ifPDFTeX
\else
\babelfont{rm}[]{Tinos}
\fi
% get rid of language-specific shorthands (see #6817):
\let\LanguageShortHands\languageshorthands
\def\languageshorthands#1{}
\ifLuaTeX
  \usepackage{selnolig}  % disable illegal ligatures
\fi
\IfFileExists{bookmark.sty}{\usepackage{bookmark}}{\usepackage{hyperref}}
\IfFileExists{xurl.sty}{\usepackage{xurl}}{} % add URL line breaks if available
\urlstyle{same}
\hypersetup{
  pdftitle={NCFV 高效变分重构、稀疏求解与 BDF2 双时间推进},
  pdflang={zh-CN},
  hidelinks,
  pdfcreator={LaTeX via pandoc}}

\title{\zihao{-3} NCFV 高效变分重构、稀疏求解与 BDF2 双时间推进}
\author{}
\date{}

\begin{document}
\maketitle

\hypertarget{ux76eeux7684ux7b26ux53f7ux4e0eux501fux9274ux8303ux56f4}{%
\subsection{1.
目的、符号与借鉴范围}\label{ux76eeux7684ux7b26ux53f7ux4e0eux501fux9274ux8303ux56f4}}

本文把王乾博士论文第 4 章的界面跳跃积分（interfacial jump
integration，IJI）思想，改写到三阶节点中心有限体积（NCFV）的对偶控制体上，并进一步给出稀疏系统的构造、可解性条件，以及
BDF2、LU-SGS 和 GMRES 的隐式推进过程。

论文中的基本做法是：每个原始单元上用零均值多项式表示重构，在所有单元界面上惩罚两侧多项式及其法向导数的跳跃，然后对全域泛函取极小值。下文保留``界面信息的加权跳跃平方和''这一变分主线，但针对
NCFV
的节点图和高效微分积分作离散调整：在每条原始节点边的中点比较两端对偶控制体多项式的值、一阶导数和二阶导数。该中点
jet 泛函与论文的面上 IJI
是相近的变分构造，二者的采样位置和导数范数不同，不能说成同一个离散格式。

以下推导采用三阶 NCFV 的二次重构。对偶控制体区域记为
\(\Omega_i\)，体积为 \(V_i=|\Omega_i|\)，其守恒变量平均值为 \begin{equation}
\bar{\boldsymbol U}_i=\frac{1}{V_i}\int_{\Omega_i}\boldsymbol U(\boldsymbol x)\,dV.\end{equation} 节点边 \(e=(i,j)\) 的几何长度为
\(\ell_e=\|\boldsymbol x_j-\boldsymbol x_i\|\)，中点为
\(\boldsymbol m_e=(\boldsymbol x_i+\boldsymbol x_j)/2\)。\(d\)
为空间维数，守恒变量个数记为 \(n_v\)。

\hypertarget{ux96f6ux5747ux503cux4e8cux6b21ux591aux9879ux5f0f}{%
\subsection{2.
零均值二次多项式}\label{ux96f6ux5747ux503cux4e8cux6b21ux591aux9879ux5f0f}}

\hypertarget{ux57faux51fdux6570ux4e0eux5e73ux5747ux7ea6ux675f}{%
\subsubsection{2.1
基函数与平均约束}\label{ux57faux51fdux6570ux4e0eux5e73ux5747ux7ea6ux675f}}

二次多项式去掉常数项后的系数数目为 \begin{equation}
M=d+\frac{d(d+1)}2=
\begin{cases}
5,&d=2,\\
9,&d=3.
\end{cases}\end{equation} 为改善拉伸网格上的条件数，对每个节点用分方向长度矩阵 \begin{equation}
H_i=\operatorname{diag}(h_{i,1},\ldots,h_{i,d}),\qquad
\boldsymbol\xi_i=H_i^{-1}(\boldsymbol x-\boldsymbol x_i)\end{equation} 做无量纲化。可取 \(h_{i,\alpha}\) 为节点对偶控制体在第 \(\alpha\)
个方向上的半跨度。三维基函数例如 \begin{equation}
\boldsymbol\psi_i=
\left[\xi,\eta,\zeta,
\tfrac12\xi^2,\xi\eta,\xi\zeta,
\tfrac12\eta^2,\eta\zeta,\tfrac12\zeta^2\right]^T,\end{equation} 二维时删去含 \(\zeta\) 的项。定义零均值基 \begin{equation}
\boldsymbol\phi_i(\boldsymbol x)=
\boldsymbol\psi_i(\boldsymbol x)-
\frac{1}{V_i}\int_{\Omega_i}\boldsymbol\psi_i(\boldsymbol x)\,dV.\end{equation} 于是
\(\int_{\Omega_i}\boldsymbol\phi_i\,dV=\boldsymbol 0\)。二维基函数的明确形式为
\begin{equation}
\boldsymbol\psi_i=[\xi,\eta,\tfrac12\xi^2,\xi\eta,\tfrac12\eta^2]^T.\end{equation} 构造零均值基时只需对这些单项式积分。令 \begin{equation}
\boldsymbol M_i^{(1)}=\int_{\Omega_i}(\boldsymbol x-\boldsymbol x_i)\,dV,\qquad
\boldsymbol M_i^{(2)}=\int_{\Omega_i}(\boldsymbol x-\boldsymbol x_i)
(\boldsymbol x-\boldsymbol x_i)^T\,dV,\end{equation} 则 \begin{equation}
\langle\boldsymbol\xi_i\rangle=H_i^{-1}\boldsymbol M_i^{(1)}/V_i,\qquad
\langle\boldsymbol\xi_i\boldsymbol\xi_i^T\rangle=
H_i^{-1}\boldsymbol M_i^{(2)}H_i^{-1}/V_i.\end{equation}
因此一次项和二次项的均值可由对偶微单纯形的解析几何矩得到，无需在重构初始化时布置
Gauss 点。节点 \(i\) 的重构为 \begin{equation}
\boxed{\boldsymbol P_i(\boldsymbol x)=\bar{\boldsymbol U}_i+
\boldsymbol C_i^T\boldsymbol\phi_i(\boldsymbol x)},\end{equation} 其中 \(\boldsymbol C_i\in\mathbb R^{M\times n_v}\)
是待定系数。零均值性质立即给出 \begin{equation}
\frac{1}{V_i}\int_{\Omega_i}\boldsymbol P_i\,dV=\bar{\boldsymbol U}_i,\end{equation} 因此重构不会改变有限体积未知量的定义，也不会破坏控制体平均守恒。

\hypertarget{ux8282ux70b9ux503cux68afux5ea6ux4e0e-hessian}{%
\subsubsection{2.2 节点值、梯度与
Hessian}\label{ux8282ux70b9ux503cux68afux5ea6ux4e0e-hessian}}

在原始节点位置恢复点值与梯度： \begin{equation}
\boldsymbol U_i^{\mathrm{node}}=\boldsymbol P_i(\boldsymbol x_i),\qquad
\boldsymbol G_i=\nabla\boldsymbol P_i(\boldsymbol x_i).\end{equation} 二阶导数 \(\nabla^2\boldsymbol P_i\)
在二次多项式内为常数。它参与下文的界面变分泛函，使系数系统能辨别二次模态；高效微分通量阶段仍可只传递
\(\boldsymbol U_i^{\mathrm{node}}\) 与
\(\boldsymbol G_i\)，不必在每个微面或 Gauss 点存 Hessian。

\hypertarget{ux4eceux754cux9762ux8df3ux8dc3ux5230ux8282ux70b9ux8fb9-jet-ux6cdbux51fd}{%
\subsection{3. 从界面跳跃到节点边 jet
泛函}\label{ux4eceux754cux9762ux8df3ux8dc3ux5230ux8282ux70b9ux8fb9-jet-ux6cdbux51fd}}

\hypertarget{ux8bbaux6587ux4e2dux7684-iji-ux5f62ux5f0f}{%
\subsubsection{3.1 论文中的 IJI
形式}\label{ux8bbaux6587ux4e2dux7684-iji-ux5f62ux5f0f}}

对论文中的单元界面 \(f\)，两侧重构分别为 \(P_L,P_R\)，中心距离为
\(d_{LR}\)。其 IJI 可概括为 \begin{equation}
I_f=\frac1{d_{LR}}\int_f\sum_{p=0}^{k}
\frac{d_{LR}^{2p}}{(p!)^2}
\left(\partial_n^p P_L-\partial_n^p P_R\right)^2dS,
\qquad I=\sum_f I_f.\end{equation} 这里 \(p=0\) 是函数值跳跃，\(p=1,2,\ldots,k\)
是法向导数跳跃；距离幂使不同导数阶的项具有一致量纲。对重构系数求驻值便得到全场稀疏线性系统，这对应论文式（4-5）---（4-11）\protect\hyperlink{ref_001}{1}。论文还证明了：在全域网格连通、重构空间满足相应非奇异性条件时，该系统对称正定\protect\hyperlink{ref_001}{1}（第4.1.2节）。

\hypertarget{ux9002ux7528ux4e8e-ncfv-ux9ad8ux6548ux91cdux6784ux7684ux79bbux6563ux754cux9762ux4fe1ux606f}{%
\subsubsection{3.2 适用于 NCFV
高效重构的离散界面信息}\label{ux9002ux7528ux4e8e-ncfv-ux9ad8ux6548ux91cdux6784ux7684ux79bbux6563ux754cux9762ux4fe1ux606f}}

NCFV
的节点对偶界面由原始边及其周围微面组成。为避免在变分构造中再次引入面上高斯点，可在每条节点边的中点比较两侧重构的完整二次
jet。对边 \(e=(i,j)\) 定义 \begin{equation}
\mathcal J_e=
\begin{bmatrix}
P_i(\boldsymbol m_e)-P_j(\boldsymbol m_e)\\[1mm]
\dfrac{\ell_e}{2}\left(\nabla P_i(\boldsymbol m_e)-\nabla P_j(\boldsymbol m_e)\right)\\[1mm]
\dfrac{\ell_e^2}{4}\left(\nabla^2P_i-\nabla^2P_j\right)
\end{bmatrix}.\end{equation} 向量梯度按 Cartesian 分量排列；Hessian 按全部 \(d^2\)
个有序分量排列，因此混合导数在 Frobenius 范数中出现两次。可定义 \begin{equation}
\boxed{\mathcal I(\boldsymbol C)=\frac12\sum_{e=(i,j)}\omega_e
\|\mathcal J_e\|_2^2,\qquad \omega_e>0.}\end{equation} 当前工作区 NCFV 变分行的尺度对应 \(\omega_e=1\)、值行权重为
1、梯度行乘 \(\ell_e/2\)、Hessian 行乘
\(\ell_e^2/4\)。若改用面面积权、不同的各阶权函数或论文的纯法向导数内积，仍可得到变分格式，但矩阵条件数和具体重构结果会改变，需重新验证。

这个泛函的含义是：在每条节点边的代表位置，使两侧的函数值、斜率和曲率尽量一致。它以节点邻接图为耦合图，符合
NCFV 的紧致数据结构。函数值项由已知的两个控制体平均值提供偏置，梯度和
Hessian 项则约束高阶系数。

\hypertarget{ux77e9ux9635ux884cux8868ux793a}{%
\subsubsection{3.3 矩阵行表示}\label{ux77e9ux9635ux884cux8868ux793a}}

令 \(q=1+d+d^2\)，定义
\(\boldsymbol e_0=[1,0,\ldots,0]^T\in\mathbb R^q\)。对边
\(e=(i,j)\)，定义矩阵 \(D_i^e\in\mathbb R^{q\times M}\)： \begin{equation}
D_i^e=
\begin{bmatrix}
\boldsymbol\phi_i(\boldsymbol m_e)^T\\
\dfrac{\ell_e}{2}\nabla\boldsymbol\phi_i(\boldsymbol m_e)^T\\
\dfrac{\ell_e^2}{4}\nabla^2\boldsymbol\phi_i^T
\end{bmatrix},\end{equation} 其中梯度与 Hessian 的每个 Cartesian
分量各占一行。对单个标量变量，界面 jet 残差写成 \begin{equation}
\boldsymbol r_e=\boldsymbol e_0(\bar u_i-\bar u_j)
+D_i^e\boldsymbol c_i-D_j^e\boldsymbol c_j.\end{equation} 向量守恒状态对每个分量使用同一套
\(D_i^e\)，只需把标量系数列向量替换为矩阵 \(\boldsymbol C_i\)。

\hypertarget{eulerlagrange-ux65b9ux7a0bux4e0eux7a00ux758fux5757ux77e9ux9635}{%
\subsection{4. Euler--Lagrange
方程与稀疏块矩阵}\label{eulerlagrange-ux65b9ux7a0bux4e0eux7a00ux758fux5757ux77e9ux9635}}

\hypertarget{ux5bf9ux7cfbux6570ux53d6ux4e00ux9636ux53d8ux5206}{%
\subsubsection{4.1
对系数取一阶变分}\label{ux5bf9ux7cfbux6570ux53d6ux4e00ux9636ux53d8ux5206}}

先考虑一个守恒分量。泛函为 \begin{equation}
\mathcal I(\boldsymbol c)=\frac12\sum_{e=(i,j)}\omega_e
\boldsymbol r_e^T\boldsymbol r_e.\end{equation} 对节点 \(i\) 的系数 \(\boldsymbol c_i\) 求导。每条相邻边向梯度贡献 \begin{equation}
(D_i^e)^T\omega_e\boldsymbol r_e,\end{equation} 全部相邻边贡献之和置零后，将含
\(\boldsymbol c_i\)、邻点系数和已知均值分别整理，得到 \begin{equation}
\boxed{
K_i\boldsymbol c_i-
\sum_{j\in N(i)}C_{ij}\boldsymbol c_j=
\sum_{j\in N(i)}\boldsymbol g_{ij}(\bar u_j-\bar u_i),}\end{equation} 其中 \begin{equation}
K_i=\sum_{e\ni i}\omega_e(D_i^e)^TD_i^e,\quad
C_{ij}=\omega_{ij}(D_i^{ij})^TD_j^{ij},\quad
\boldsymbol g_{ij}=\omega_{ij}(D_i^{ij})^T\boldsymbol e_0.\end{equation} 若边权在不同边上变化，以上求和中的 \(\omega_e\)
按边分别取值。对所有节点组装得到 \begin{equation}
\boxed{\mathbb A\boldsymbol c=\boldsymbol b,\qquad
\mathbb A_{ii}=K_i,\quad
\mathbb A_{ij}=-C_{ij}\ (j\in N(i)),\quad
\mathbb A_{ij}=0\ (j\notin N(i)).}\end{equation} 右端为 \begin{equation}
\boldsymbol b_i=\sum_{j\in N(i)}\boldsymbol g_{ij}(\bar u_j-\bar u_i).\end{equation} 因为反向边给出 \(C_{ji}=C_{ij}^T\)，全局矩阵 \(\mathbb A\)
对称。它的块稀疏图就是原始节点边图：每个节点只有一个 \(M\times M\)
对角块以及相邻节点的非对角块。因此储存和矩阵向量乘法的复杂度随节点数和边数线性增长，块内的小矩阵尺寸固定为
5 或 9。

将全部守恒变量并列，方程成为 \begin{equation}
\mathbb A\boldsymbol C=\boldsymbol B,\end{equation} 其中 \(\boldsymbol C_i\in\mathbb R^{M\times n_v}\)，同一几何矩阵对
\(n_v\) 个分量重复求解。这允许复用分解、预条件器和稀疏矩阵结构。

\hypertarget{ux5c40ux90e8ux56faux5b9aux70b9-block-gaussseidel}{%
\subsubsection{4.2 局部固定点 / Block
Gauss--Seidel}\label{ux5c40ux90e8ux56faux5b9aux70b9-block-gaussseidel}}

若逐节点解局部对角块，可写为 \begin{equation}
\boldsymbol c_i^{(s+1)}=K_i^{-1}\left[
\sum_{j<i}C_{ij}\boldsymbol c_j^{(s+1)}+
\sum_{j>i}C_{ij}\boldsymbol c_j^{(s)}+
\boldsymbol b_i\right].\end{equation} 这是按节点顺序的 block Gauss--Seidel。带松弛的 block SOR 为 \begin{equation}
\boldsymbol c_i^{(s+1)}\leftarrow
(1-\omega)\boldsymbol c_i^{(s)}+
\omega K_i^{-1}\left[
\sum_{j<i}C_{ij}\boldsymbol c_j^{(s+1)}+
\sum_{j>i}C_{ij}\boldsymbol c_j^{(s)}+\boldsymbol b_i\right].\end{equation} 当前代码所做的``对每个节点直接解
\(K_i\)，再重复邻点系数更新''，外层轮数是全局耦合迭代次数，并不是反复近似局部
\(K_i\) 的逆。使用最小二乘系数作初值可减少外层误差；每轮 MPI ghost
同步后，各 rank 的并行更新不再等同于全局串行排序的精确
Gauss--Seidel，实际收敛率需用重构残差或二次场固定点误差监控。

\hypertarget{ux6709ux89e3ux6027ux552fux4e00ux6027ux4e0eux7cbeux786eux6027}{%
\subsection{5.
有解性、唯一性与精确性}\label{ux6709ux89e3ux6027ux552fux4e00ux6027ux4e0eux7cbeux786eux6027}}

\hypertarget{ux534aux6b63ux5b9aux6027}{%
\subsubsection{5.1 半正定性}\label{ux534aux6b63ux5b9aux6027}}

由泛函 Hessian 的构造，任意系数扰动 \(\boldsymbol z\) 满足 \begin{equation}
\boldsymbol z^T\mathbb A\boldsymbol z=
\sum_{e=(i,j)}\omega_e
\|D_i^e\boldsymbol z_i-D_j^e\boldsymbol z_j\|_2^2\ge 0.\end{equation} 因此 \(\mathbb A\)
总是对称半正定。这里已知平均值只进入右端，不改变系统矩阵。

\hypertarget{ux4f55ux65f6ux4e25ux683cux6b63ux5b9a}{%
\subsubsection{5.2
何时严格正定}\label{ux4f55ux65f6ux4e25ux683cux6b63ux5b9a}}

半正定要成为正定，还需排除零能量的非零系数场。一个直接的充分条件是：

\begin{enumerate}
\def\labelenumi{\arabic{enumi}.}
\tightlist
\item
  每个连通分量上的节点图连通，所有边权严格为正，边的坐标均在两侧多项式的同一个几何坐标框架中比较（周期边需采用一致的周期展开位移）；
\item
  每个节点的局部行矩阵 \(\{D_i^e:e\ni i\}\) 满列秩，使 \(K_i\)
  正定、局部更新有定义；
\item
  每个连通分量上，二次多项式空间的控制体平均泛函具有满列秩。二维需至少有
  6 个几何独立的控制体平均约束，三维需至少 10
  个；数量足够仍不够，还必须几何非退化。
\end{enumerate}

证明如下。若
\(\boldsymbol z^T\mathbb A\boldsymbol z=0\)，每条边的值、梯度和 Hessian
jet 跳跃都为零。由于局部多项式次数不超过 2，Hessian
为常数；沿连通节点图传播，所有单元扰动多项式具有相同
Hessian。中点梯度相等继而使线性项相同，中点值相等再使常数项相同。因此所有局部扰动多项式是同一个全局二次多项式
\(q(\boldsymbol x)\)。每个扰动多项式在自己的对偶控制体上均值为零，故 \begin{equation}
\int_{\Omega_i}q(\boldsymbol x)\,dV=0\quad\text{对该连通分量的所有 }i.\end{equation} 若控制体平均对 \(\mathbb P_2\) 满列秩，则只能有
\(q\equiv0\)，从而所有 \(\boldsymbol z_i=0\)。因此 \(\mathbb A\)
对称正定，方程有唯一解。

这也指出了可解性论证的边界：论文的 SPD 结论不能脱离网格非退化和 k-exact
/ 均值泛函满秩条件直接照搬到任意节点图。局部 \(K_i\) 可逆只说明每个
block
更新可算，不能单独证明全局矩阵无零空间。若条件不满足，应扩充节点模板、加入有物理依据的边界/周期约束，或用秩揭示
QR/SVD 检查并处理零空间；简单加 \(\epsilon I\)
虽可强制可逆，却会改变重构并破坏严格多项式精确性。

\hypertarget{ux4e8cux6b21ux591aux9879ux5f0fux7cbeux786eux6027}{%
\subsubsection{5.3
二次多项式精确性}\label{ux4e8cux6b21ux591aux9879ux5f0fux7cbeux786eux6027}}

设精确场 \(u(\boldsymbol x)\in\mathbb P_2\)，且输入 \(\bar u_i\) 是其在
\(\Omega_i\) 上的真实平均。零均值基可精确表示
\(u-\bar u_i\)，所以存在一组系数使每个 \(P_i=u\)。这时每条边 jet
残差为零，\(\mathcal I=0\)。若上述唯一性条件成立，该零泛函解就是唯一极小值解，故重构对任意二次多项式精确。相应的一阶导数在光滑场上具备构造三阶
NCFV 所需的二阶精度基础。

三阶收敛仍需形状规则网格、稳定重构、足够收敛的变分迭代、适当的边界处理和未触发降阶的限制器；Roe
耗散、非线性物理通量的高效积分近似及时间误差也会影响观测阶数。不能仅凭二次精确性宣称所有流动情形均达到三阶。

\hypertarget{ux53d8ux5206ux7cfbux7edfux7684ux6570ux503cux89e3ux6cd5ux4e0eux6548ux7387}{%
\subsection{6.
变分系统的数值解法与效率}\label{ux53d8ux5206ux7cfbux7edfux7684ux6570ux503cux89e3ux6cd5ux4e0eux6548ux7387}}

\hypertarget{ux53efux9009ux6c42ux89e3ux5668}{%
\subsubsection{6.1 可选求解器}\label{ux53efux9009ux6c42ux89e3ux5668}}

在严格正定条件下，建议按问题规模选择：

\begin{itemize}
\tightlist
\item
  \textbf{小网格直接法}：稀疏 Cholesky，用于基准计算和验证唯一解；
\item
  \textbf{PCG}：\(\mathbb A\) 对称正定时的首选迭代法，适合用 block
  Jacobi \(\operatorname{diag}(K_i)\)、不完全 Cholesky 或 AMG 预条件；
\item
  \textbf{Block Gauss--Seidel / SOR}：重用每个节点的小型 \(K_i\)
  分解，内存少，适合与时间推进耦合；
\item
  \textbf{GMRES}：适用于因近似线性化、非对称边界项或非对称预条件而形成的非对称变分方程，不应仅因``GMRES
  通用''就替代 SPD 系统上的 PCG。
\end{itemize}

可用边项直接做 matrix-free 乘法。给定向量 \(\boldsymbol z\)，先逐边计算
\(\boldsymbol q_e=D_i^e\boldsymbol z_i-D_j^e\boldsymbol z_j\)，再把
\((D_i^e)^T\omega_e\boldsymbol q_e\) 和
\(-(D_j^e)^T\omega_e\boldsymbol q_e\) 加回两端。这样不用存储全局 CSR
矩阵；但每轮仍需邻点交换与一次全局内积归约（PCG）。

\hypertarget{ux4e0eux9ad8ux6548-differential-ux901aux91cfux7684ux8854ux63a5}{%
\subsubsection{6.2 与高效 differential
通量的衔接}\label{ux4e0eux9ad8ux6548-differential-ux901aux91cfux7684ux8854ux63a5}}

每次残差计算先由当前控制体平均 \(\bar{\boldsymbol U}\)
解变分系数，随后只抽取 \begin{equation}
\boldsymbol U_i^{\mathrm{node}}=\boldsymbol P_i(\boldsymbol x_i),\qquad
\boldsymbol G_i=\nabla\boldsymbol P_i(\boldsymbol x_i).\end{equation}
把这两个场送入预计算的对偶体积/宏面微分权重。对光滑状态，物理通量梯度按链式法则计算：
\begin{equation}
\partial_\alpha\boldsymbol F_\beta(\boldsymbol U_i)
=\frac{\partial\boldsymbol F_\beta}{\partial\boldsymbol U}
\bigg|_{\boldsymbol U_i}\partial_\alpha\boldsymbol U_i.\end{equation} 于是体积积分、宏面状态平均和中心物理通量积分使用点值与一阶导数的压缩
stencil，不需要在时间步内重建 Gauss
点；每条原始边对应的对偶宏面仍只做一次 Roe/Riemann
求解，端点残差用相反符号累加以保持守恒。变分重构阶段使用
Hessian，是为了确定二次系数；它不意味着高效 flux kernel 必须读取
Hessian。

效率上的代价也应明确：与只保存梯度的 LS
路径相比，变分重构必须在系数迭代中保留每节点 \(M\times n_v\)
个二次系数，并进行邻点同步。可以在残差阶段只保留恢复出的节点值和梯度，但不能在变分系数尚未收敛时丢弃系数。几何矩阵
\(K_i\)
和变分算子只依赖网格，可预计算；系数解可用前一残差/时间步作初值。

如果希望每个时间残差都把变分方程解到很紧的容差，成本可能超过 flux
积分节省的成本。可使用 warm start 和固定的少量 block
sweeps，并把``重构迭代误差''纳入总非线性缺陷监控。若 GMRES 需要严格的
Newton--Krylov
线性化，则矩阵向量乘法应包括重构映射的响应；冻结变分系数得到的是近似
Jacobian / Picard 预条件，需在外层更新重构并验证缺陷收敛。

\hypertarget{ux5b88ux6052ux534aux79bbux6563ux65b9ux7a0b}{%
\subsection{7.
守恒半离散方程}\label{ux5b88ux6052ux534aux79bbux6563ux65b9ux7a0b}}

对静止网格，用节点对偶体平均 \(\bar{\boldsymbol U}_i\)
作未知量，记广延守恒量 \begin{equation}
\boldsymbol Q_i=V_i\bar{\boldsymbol U}_i.\end{equation} 内部宏面 \(e\) 的有向积分通量记为
\(\widehat{\boldsymbol I}_e\)，节点关联符号为
\(B_{ie}=\pm1\)。有限体积残差为 \begin{equation}
\boldsymbol{\mathcal R}_i(\bar{\boldsymbol U})=
-\sum_{e\in E(i)}B_{ie}\widehat{\boldsymbol I}_e-
\boldsymbol I_{\partial V_i},\qquad
\boxed{\frac{d\boldsymbol Q_i}{dt}=\boldsymbol{\mathcal R}_i(\bar{\boldsymbol U}).}\end{equation} 同一内部面通量在两端符号相反，因此全域内部通量严格抵消。定义
\(J_{ij}=\partial\boldsymbol{\mathcal R}_i/\partial\bar{\boldsymbol U}_j\)。精确变分重构映射会让
(J)
的数学依赖通过全局重构系统扩展到整个连通网格；工程上常用冻结重构的局部面
Jacobian 作为近似，或用矩阵自由差分/自动微分把重构影响计入
Jacobian-vector product。

\hypertarget{bdf2-ux7269ux7406ux65f6ux95f4ux79bbux6563}{%
\subsection{8. BDF2
物理时间离散}\label{bdf2-ux7269ux7406ux65f6ux95f4ux79bbux6563}}

\hypertarget{ux5b9aux6b65ux957f-bdf2}{%
\subsubsection{8.1 定步长 BDF2}\label{ux5b9aux6b65ux957f-bdf2}}

对于均匀物理步长 \(\Delta t\)，在新时刻隐式离散为 \begin{equation}
\boxed{
\frac{3\boldsymbol Q_i^{n+1}-4\boldsymbol Q_i^n+
\boldsymbol Q_i^{n-1}}{2\Delta t}
=\boldsymbol{\mathcal R}_i(\bar{\boldsymbol U}^{n+1}).}\end{equation} 因 \(V_i\) 固定，也可逐控制体写成平均状态形式 \begin{equation}
V_i\frac{3\bar{\boldsymbol U}_i^{n+1}-4\bar{\boldsymbol U}_i^n+
\bar{\boldsymbol U}_i^{n-1}}{2\Delta t}
=\boldsymbol{\mathcal R}_i(\bar{\boldsymbol U}^{n+1}).\end{equation} 其时间导数近似为二阶精度。第一物理步没有 \(n-1\) 层，使用后向 Euler
启动： \begin{equation}
\frac{\boldsymbol Q_i^1-\boldsymbol Q_i^0}{\Delta t}
=\boldsymbol{\mathcal R}_i(\bar{\boldsymbol U}^{1}).\end{equation} 启动步只用于建立两层历史；之后每步保持 \(\bar{\boldsymbol U}^n\) 和
\(\bar{\boldsymbol U}^{n-1}\)
不变，直至新物理层内迭代收敛，再整体移动历史。不能在伪时间子迭代中提前覆盖历史场。

变步长时令 \(r=\Delta t_n/\Delta t_{n-1}\)，导数公式为 \begin{equation}
\frac1{\Delta t_n}\left[
\frac{1+2r}{1+r}\boldsymbol Q^{n+1}
-(1+r)\boldsymbol Q^n+\frac{r^2}{1+r}\boldsymbol Q^{n-1}\right].\end{equation} 定步长格式对应
\(r=1\)。时间步变化时必须同时更新残差和隐式对角项中的首系数，不能仍使用
\(3/2\)。

\hypertarget{ux53ccux65f6ux95f4ux4f2aux65f6ux95f4ux5185ux8fedux4ee3}{%
\subsubsection{8.2
双时间伪时间内迭代}\label{ux53ccux65f6ux95f4ux4f2aux65f6ux95f4ux5185ux8fedux4ee3}}

为把每个非线性 BDF2 方程变成可迭代的伪稳态问题，在新物理步内引入伪时间
\(\tau\)： \begin{equation}
V_i\frac{\partial\bar{\boldsymbol U}_i}{\partial\tau}
+V_i\frac{3\bar{\boldsymbol U}_i-4\bar{\boldsymbol U}_i^n+
\bar{\boldsymbol U}_i^{n-1}}{2\Delta t}
=\boldsymbol{\mathcal R}_i(\bar{\boldsymbol U}).\end{equation} 伪时间导数收敛为零后，所得解正是原 BDF2 离散解。用伪时间后向 Euler
修正，从当前内迭代状态 \(\bar{\boldsymbol U}^{(m)}\) 求增量
\(\delta\bar{\boldsymbol U}\)。定义物理缺陷 \begin{equation}
\boldsymbol F_i^{(m)}=
V_i\frac{3\bar{\boldsymbol U}_i^{(m)}-4\bar{\boldsymbol U}_i^n+
\bar{\boldsymbol U}_i^{n-1}}{2\Delta t}
-\boldsymbol{\mathcal R}_i(\bar{\boldsymbol U}^{(m)}).\end{equation} 一阶线性化后求 \begin{equation}
\boxed{
\sum_j\left[
\left(\frac{V_i}{\Delta\tau_i}+\frac{3V_i}{2\Delta t}\right)
\boldsymbol I\,\delta_{ij}-J_{ij}
\right]\delta\bar{\boldsymbol U}_j=-\boldsymbol F_i^{(m)},\qquad
\bar{\boldsymbol U}^{(m+1)}=
\bar{\boldsymbol U}^{(m)}+\lambda\delta\bar{\boldsymbol U}.}\end{equation} \(\Delta\tau_i\) 可按局部 CFL 选取，影响伪时间收敛速度；物理步长
\(\Delta t\)
决定真实时间离散精度。每个物理步应检查物理缺陷范数，而不只检查增量范数。对于守恒量方程，BDF2
的物理质量项作用于全部守恒变量；若用于压力不带物理时间导数的人工压缩系统，则需另设物理质量矩阵，不能直接套用此处全变量单位矩阵。

\hypertarget{lu-sgs-ux7ebfux6027ux6c42ux89e3ux4e0e-bdf2-ux6821ux6b63}{%
\subsection{9. LU-SGS 线性求解与 BDF2
校正}\label{lu-sgs-ux7ebfux6027ux6c42ux89e3ux4e0e-bdf2-ux6821ux6b63}}

\hypertarget{ux77e9ux9635ux5206ux88c2}{%
\subsubsection{9.1 矩阵分裂}\label{ux77e9ux9635ux5206ux88c2}}

将双时间线性系统记为 \begin{equation}
\mathbb A\,\delta\boldsymbol U=\boldsymbol b,
\qquad \mathbb A=\mathbb D+\mathbb L+\mathbb U,\end{equation} 其中 \(\mathbb D\) 是节点对角块，\(\mathbb L\)、\(\mathbb U\)
分别为按节点编号排序后的严格下、上三角耦合块。精确 Jacobian
可能昂贵，可用局部 Roe/Rusanov 波速构造对角占优近似，例如 \begin{equation}
D_i\approx
\left[\frac{V_i}{\Delta\tau_i}+\frac{3V_i}{2\Delta t}
+\frac12\sum_{e\ni i}S_e\lambda_e\right]\boldsymbol I,\end{equation} 其中 \(S_e\) 是宏面面积，\(\lambda_e\)
为正的对流谱半径；粘性问题可再加入粘性谱半径估计。面间块由冻结重构后数值通量的局部导数给出。该式是预条件近似，不是完整高阶离散的精确
Jacobian。

\hypertarget{ux5bf9ux79f0ux524dux540eux626bux63cf}{%
\subsubsection{9.2
对称前后扫描}\label{ux5bf9ux79f0ux524dux540eux626bux63cf}}

LU-SGS 用如下因子近似 \begin{equation}
\mathbb P=(\mathbb D+\mathbb L)\mathbb D^{-1}(\mathbb D+\mathbb U),\end{equation} 对每个增量方程先正向再反向扫描。

\textbf{正向扫描}，按节点从小到大： \begin{equation}
(D_i+L_{ii})\boldsymbol y_i=
\boldsymbol b_i-\sum_{j<i}L_{ij}\boldsymbol y_j.\end{equation} 实际 \(L_{ii}=0\)，故每步只需解一个 \(n_v\times n_v\) 的局部对角块。

\textbf{反向扫描}，按节点从大到小： \begin{equation}
D_i\delta\boldsymbol U_i=
D_i\boldsymbol y_i-\sum_{j>i}U_{ij}\delta\boldsymbol U_j,\end{equation} 即 \begin{equation}
\delta\boldsymbol U_i=\boldsymbol y_i-
D_i^{-1}\sum_{j>i}U_{ij}\delta\boldsymbol U_j.\end{equation} 所得 \(\delta\boldsymbol U\) 是一次 LU-SGS
近似解；若只做一组前后扫，它是低成本近似线性求解，若再以线性残差做
correction sweeps，则继续逼近原线性系统。伪代码为：

\begin{Shaded}
\begin{Highlighting}[]
\NormalTok{输入：当前状态 U、物理历史 U\^{}n/U\^{}(n{-}1)、缺陷 b、近似块 L/D/U}
\NormalTok{y = 0}
\NormalTok{for i = 1,...,N:}
\NormalTok{    y\_i = D\_i\^{}\{{-}1\} (b\_i {-} sum\_\{j\textless{}i\} L\_ij y\_j)}
\NormalTok{deltaU = 0}
\NormalTok{for i = N,...,1:}
\NormalTok{    deltaU\_i = y\_i {-} D\_i\^{}\{{-}1\} sum\_\{j\textgreater{}i\} U\_ij deltaU\_j}
\NormalTok{U \textless{}{-} U + relaxation * deltaU}
\NormalTok{若物理缺陷未收敛，则重新计算重构、通量和下一次修正}
\end{Highlighting}
\end{Shaded}

节点顺序和跨 MPI 分区通信会影响 LU-SGS 预条件质量。分布式实现通常在
forward/backward sweep 之间同步 ghost；跨 rank
的节点只能使用当轮收到的边界值，因而是近似 SGS。通过图着色、多色 SGS
或在 rank 间采用 block Jacobi
可改善并行实现，但都应以物理缺陷收敛为判断依据。

\hypertarget{gmres-ux4e0e-lu-sgs-ux9884ux6761ux4ef6}{%
\subsection{10. GMRES 与 LU-SGS
预条件}\label{gmres-ux4e0e-lu-sgs-ux9884ux6761ux4ef6}}

\hypertarget{ux4e3aux4ec0ux4e48ux7528ux4e8eux6d41ux52a8ux65f6ux95f4ux65b9ux7a0b}{%
\subsubsection{10.1
为什么用于流动时间方程}\label{ux4e3aux4ec0ux4e48ux7528ux4e8eux6d41ux52a8ux65f6ux95f4ux65b9ux7a0b}}

BDF2 线性系统含有对流通量
Jacobian，通常非对称；中心/耗散项、边界条件、粘性项也会影响谱性质。GMRES
不要求矩阵对称正定，能最小化 Krylov 子空间中的线性残差。LU-SGS
可作为左预条件器，形成 \begin{equation}
\mathbb P^{-1}\mathbb A\,\delta\boldsymbol U=
\mathbb P^{-1}\boldsymbol b.\end{equation} 一次 \(\mathbb P^{-1}\) 应用就是上一节的一组 LU-SGS
前后扫描。若使用会随 Krylov 步变化的非线性预条件操作，则应换用 flexible
GMRES；若每次都是固定的线性 LU-SGS 近似，标准左预条件 GMRES 即可。

\hypertarget{arnoldigmres-ux6b65ux9aa4}{%
\subsubsection{10.2 Arnoldi--GMRES
步骤}\label{arnoldigmres-ux6b65ux9aa4}}

令 \(\mathbb M=\mathbb P^{-1}\)，从 \(\delta\boldsymbol U_0\) 开始：

\begin{enumerate}
\def\labelenumi{\arabic{enumi}.}
\tightlist
\item
  计算左预条件初始残差
  \(\boldsymbol r_0=\mathbb M(\boldsymbol b-\mathbb A\delta\boldsymbol U_0)\)，令
  \(\beta=\|\boldsymbol r_0\|_2\)、\(\boldsymbol v_1=\boldsymbol r_0/\beta\)。分布式范数含一次
  MPI 全局求和。
\item
  对 \(j=1,\ldots,m\)：计算
  \(\boldsymbol w=\mathbb M\mathbb A\boldsymbol v_j\)，用 modified
  Gram--Schmidt 对 \(\boldsymbol v_1,\ldots,\boldsymbol v_j\)
  正交化，得到 Hessenberg 矩阵第 \(j\) 列及新向量
  \(\boldsymbol v_{j+1}\)。
\item
  对 Hessenberg 小矩阵逐列应用 Givens 旋转，更新最小二乘残差估计；若
  \(\|r_j\|/\|r_0\|<\varepsilon_{\mathrm{lin}}\)，结束当前 Krylov 循环。
\item
  求解小型上三角系统并更新
  \(\delta\boldsymbol U=\delta\boldsymbol U_0+V_m\boldsymbol y_m\)。未收敛时以新残差重启，直到达到容差或最大重启数。
\end{enumerate}

核心伪代码：

\begin{Shaded}
\begin{Highlighting}[]
\NormalTok{deltaU = 0}
\NormalTok{repeat for each restart:}
\NormalTok{    r = M * (b {-} A * deltaU)}
\NormalTok{    beta = norm(r)}
\NormalTok{    v[1] = r / beta}
\NormalTok{    for j = 1,...,restart\_length:}
\NormalTok{        w = M * (A * v[j])       \# A 的作用可矩阵自由计算}
\NormalTok{        for k = 1,...,j:}
\NormalTok{            H[k,j] = dot(w, v[k])}
\NormalTok{            w {-}= H[k,j] * v[k]}
\NormalTok{        H[j+1,j] = norm(w)}
\NormalTok{        v[j+1] = w / H[j+1,j]}
\NormalTok{        更新 Givens 旋转和残差估计}
\NormalTok{        若达到线性容差则 break}
\NormalTok{    解 min || beta*e1 {-} H*y ||，并令 deltaU += V*y}
\NormalTok{直到线性容差或重启上限}
\end{Highlighting}
\end{Shaded}

\hypertarget{jacobianvector-product-ux4e0eux91cdux6784ux8026ux5408}{%
\subsubsection{10.3 Jacobian--vector product
与重构耦合}\label{jacobianvector-product-ux4e0eux91cdux6784ux8026ux5408}}

GMRES 只需要线性算子 \(\mathbb A\boldsymbol v\)，不必显式存储完整
Jacobian。可用解析面通量导数、自动微分，或残差差分 \begin{equation}
\mathbb A\boldsymbol v\approx
\frac{\boldsymbol F(\bar{\boldsymbol U}+\epsilon\boldsymbol v)-
\boldsymbol F(\bar{\boldsymbol U})}{\epsilon}.\end{equation} 最后一种做法要求每次 \(\boldsymbol F\)
评估都按一致规则完成变分重构。冻结系数得到的 \(\mathbb A\boldsymbol v\)
是近似线性化，可用于预条件或 inexact
Newton，但不能称为对完整变分离散残差的精确 Jacobian。

建议的双时间内层顺序是：

\begin{Shaded}
\begin{Highlighting}[]
\NormalTok{固定 U\^{}n、U\^{}(n{-}1)，令 U* = U\^{}n}
\NormalTok{for pseudo iteration m = 0,...:}
\NormalTok{    从 U* warm{-}start 并更新变分系数 C}
\NormalTok{    用 EfficientDifferential 计算积分残差 R(U*)}
\NormalTok{    形成 BDF2 物理缺陷 F(U*)}
\NormalTok{    若 ||F|| 已达标，结束当前物理步}
\NormalTok{    构造/作用 A = V/dtau + 3V/(2dt) {-} dR/dU}
\NormalTok{    用 LU{-}SGS 直接近似求解，或用 GMRES(LU{-}SGS) 求解 A*dU = {-}F}
\NormalTok{    U* \textless{}{-} U* + relaxation*dU}
\NormalTok{提交 U\^{}(n+1)=U*，之后才移动两层时间历史}
\end{Highlighting}
\end{Shaded}

\hypertarget{ux5b9eux65bdux65f6ux5e94ux62a5ux544aux7684ux6b8bux5deeux4e0eux9650ux5236}{%
\subsection{11.
实施时应报告的残差与限制}\label{ux5b9eux65bdux65f6ux5e94ux62a5ux544aux7684ux6b8bux5deeux4e0eux9650ux5236}}

建议分别监控：

\begin{enumerate}
\def\labelenumi{\arabic{enumi}.}
\tightlist
\item
  \textbf{变分重构残差}：\(\|\mathbb A C-B\|\) 或界面 jet 泛函
  \(\mathcal I\) 的相对下降；
\item
  \textbf{物理 BDF2
  缺陷}：\(\|3Q^{n+1}-4Q^n+Q^{n-1}-2\Delta t\mathcal R(Q^{n+1})\|\)
  的体积/变量归一化范数；
\item
  \textbf{线性求解残差}：\(\|b-A\delta U\|/\|b\|\)，不要只报告 GMRES
  内部的预条件残差；
\item
  \textbf{守恒误差}：全域质量、动量和能量的边界通量收支差。
\end{enumerate}

激波附近应在变分系数得到后实施适当限制，限制器可能使局部多项式不再满足未限制变分方程，也会降低局部阶数。物理步长收敛判据应基于限制后的实际残差。有限次重构
sweeps、冻结 Jacobian、LU-SGS
扫描和松弛都是求解策略；它们不会改变目标方程的形式，但若内迭代未充分收敛，会引入额外迭代误差。

\hypertarget{ux4e0eux6e90ux8bbaux6587ux53caux79fbux690dux8bf4ux660eux7684ux5bf9ux5e94}{%
\subsection{12.
与源论文及移植说明的对应}\label{ux4e0eux6e90ux8bbaux6587ux53caux79fbux690dux8bf4ux660eux7684ux5bf9ux5e94}}

\begin{itemize}
\tightlist
\item
  王乾博士论文《非结构网格紧致高精度有限体积方法》第 4 章第 4.1
  节（印刷页 74--79）给出零均值重构、面上 IJI、全局稀疏方程、SPD 证明与
  Block Jacobi / Gauss--Seidel / SOR 解法；第 2 章第 2.3.2 节（印刷页
  25--27）介绍隐式残差线性化、LU-SGS 与
  GMRES+LU-SGS\protect\hyperlink{ref_001}{1}。
\item
  本文将其``全域界面泛函的极小化''转化为节点边中点的值---梯度---Hessian
  jet 泛函，以适配二次 NCFV 重构和压缩 differential
  权重。该改写的矩阵与当前 NCFV 变分重构行算子相符，但它是 NCFV
  离散选择，不是论文逐式原样复制。
\item
  用户提供的《ncfv\_efficient\_differential\_porting\_zh.md》描述了高效算法的对偶几何、解析微分积分权重、点值恢复和宏面一次
  Riemann 求解；本文沿用这些定义，并把 BDF2
  作为隐式时间推进扩展单独推导。王乾论文中非定常示例使用的是
  SDIRK4，因此本文的 BDF2 方程不是对论文时间格式的复述。
\end{itemize}

\hypertarget{ux53c2ux8003ux6587ux732e}{%
\subsection{参考文献}\label{ux53c2ux8003ux6587ux732e}}

\protect\hypertarget{ref_001}{}{}{[}1{]} 王乾.
非结构网格紧致高精度有限体积方法{[}D{]}. 北京: 清华大学, 2017.

\end{document}
